% =============================================================================
% Patent Landscape Analysis — DRQ / SSIS
% Attorney Review Document
% UBERMASCHINE Pvt Ltd | May 2026
% Classification: ATTORNEY-CONFIDENTIAL
% =============================================================================
\documentclass[11pt,a4paper]{article}

% --- Geometry & Spacing ---
\usepackage[margin=1in]{geometry}
\usepackage{setspace}\onehalfspacing
\setlength{\parskip}{0.4em}\setlength{\parindent}{0pt}

% --- Typography ---
\usepackage[T1]{fontenc}
\usepackage{helvet}\renewcommand{\familydefault}{\sfdefault}

% --- Section Formatting ---
\usepackage{titlesec}
\titleformat{\section}{\Large\bfseries}{\thesection.}{0.6em}{}
\titleformat{\subsection}{\large\bfseries}{\thesubsection.}{0.6em}{}
\titleformat{\subsubsection}{\normalsize\bfseries}{\thesubsubsection.}{0.6em}{}

% --- Tables ---
\usepackage{booktabs}
\usepackage{tabularx}
\usepackage{array}
\newcolumntype{Y}{>{\raggedright\arraybackslash}X}

% --- Math ---
\usepackage{amsmath,amssymb}

% --- Lists ---
\usepackage{enumitem}
\setlist[itemize]{leftmargin=*,itemsep=0.1em,topsep=0.2em}
\setlist[enumerate]{leftmargin=*,itemsep=0.1em,topsep=0.2em}

% --- Colors ---
\usepackage{xcolor}
\definecolor{attgreen}{RGB}{0,128,0}
\definecolor{attred}{RGB}{180,0,0}
\definecolor{attyellow}{RGB}{180,140,0}
\definecolor{attgray}{RGB}{100,100,100}

% --- Header / Footer ---
\usepackage{fancyhdr}
\pagestyle{fancy}\fancyhf{}
\lhead{\footnotesize\textcolor{attgray}{Patent Landscape Analysis — DRQ/SSIS}}
\rhead{\footnotesize\textcolor{attgray}{UBERMASCHINE Pvt Ltd}}
\cfoot{\thepage}
\renewcommand{\headrulewidth}{0.4pt}

% =============================================================================
\begin{document}

% --- Title Page ---
\begin{center}
\vspace*{1cm}
{\LARGE\bfseries Patent Landscape Analysis}\\[0.3em]
{\large DRQ / SSIS Technology Stack}\\[1em]
{\normalsize Attorney Review Document}\\[0.3em]
{\small UBERMASCHINE Pvt Ltd}\\[0.3em]
{\small May 2026}\\[1em]
\rule{0.6\textwidth}{0.5pt}\\[0.5em]
{\small\textbf{Classification:} ATTORNEY-CONFIDENTIAL}
\end{center}

\vspace{1cm}

% =============================================================================
\section{Executive Summary}

This report provides a comprehensive patent landscape analysis for the DRQ/SSIS technology stack, conducted across WIPO PatentScope, USPTO, EPO, CNIPA, and Google Patents databases (2000--2026).

\textbf{Key Findings:}
\begin{itemize}
\item \textcolor{attgreen}{\textbf{3 of 5 core innovations are CLEAR}} — no prior art patents found
\item \textcolor{attyellow}{\textbf{1 innovation has PARTIAL OVERLAP}} — Chinese patent (CN104199022B) on general IMM for hypersonic tracking
\item \textcolor{attgreen}{\textbf{1 innovation is LARGELY CLEAR}} — generic GNSS spoofing patents exist but NavIC+SSA context is novel
\end{itemize}

\textbf{Critical Urgency:}
\begin{itemize}
\item India has \textbf{NO grace period} for patent filing after public disclosure
\item IAF ADITI submission will become public → must file provisional applications \textbf{before} submission
\item Consider Section 39 secrecy direction for defence-related inventions
\end{itemize}

% =============================================================================
\section{Core Innovation Analysis}

\subsection{Innovation 1: DPHNO — Dissipative Port-Hamiltonian Neural Operator for Orbital Propagation}

\textbf{Status: \textcolor{attgreen}{CLEAR — No competing patent found}}

\subsubsection*{Search Scope}
\begin{itemize}
\item Keywords: "port-hamiltonian neural network", "dissipative", "satellite orbit", "orbital propagation", "space situational awareness"
\item Databases: WIPO PatentScope, Google Patents, USPTO, EPO, CNIPA
\item Timeframe: 2000--2026
\end{itemize}

\subsubsection*{Findings}
\begin{itemize}
\item Zero patents combining Port-Hamiltonian neural networks with satellite orbit prediction
\item Research papers exist (arxiv 2023--2025) but \textbf{none filed as patents} in SSA context
\item The Cholesky-enforced dissipation $R = LL^\top$ applied to satellite dynamics is \textbf{genuinely novel patent territory}
\end{itemize}

\subsubsection*{Closest Prior Art}
\begin{itemize}
\item \textbf{US20140074766A1} (Numerica Corp, 2013) — uses implicit Runge-Kutta for orbit state propagation
  \begin{itemize}
  \item Overlap: None — pure numerical methods, zero neural network component
  \item Relevance: Low — different technical approach
  \end{itemize}
\end{itemize}

\subsubsection*{Patentability Assessment}
\begin{itemize}
\item \textbf{Novelty:} High — no prior art found
\item \textbf{Non-obviousness:} Strong — combining Port-Hamiltonian theory with orbital dynamics is non-obvious to those skilled in the art
\item \textbf{Industrial applicability:} Clear — direct application to space domain awareness systems
\end{itemize}

\subsubsection*{Recommended IPC Classification}
\begin{itemize}
\item Primary: G06N3/00 (Arrangements based on computational models)
\item Secondary: B64G1/10 (Spacecraft guidance)
\end{itemize}

\subsubsection*{Suggested Independent Claims}

\textbf{Claim 1 (Method):} A computer-implemented method for orbital propagation of a space object, the method comprising: obtaining a state vector of the space object; computing a time derivative of the state vector using a dissipative Port-Hamiltonian neural operator, wherein the neural operator implements: a canonical symplectic matrix $J$; a learned dissipation matrix $R$ parameterized as $R = LL^\top$, where $L$ is a Cholesky factor; a scalar Hamiltonian network $H$ parameterized by neural network weights; and a control coupling network $g$; propagating the state vector forward in time using the computed time derivative; wherein the Cholesky parameterization of $R$ guarantees positive semi-definiteness, ensuring energy dissipation in the absence of external forcing.

\textbf{Claim 2 (System):} A system for orbital propagation, comprising: one or more processors; memory storing instructions that, when executed by the processors, cause the system to: obtain a state vector of a space object; compute a time derivative of the state vector using a dissipative Port-Hamiltonian neural operator; wherein the neural operator comprises: a canonical symplectic matrix $J$; a learned dissipation matrix $R = LL^\top$, where $L$ is a Cholesky factor; a scalar Hamiltonian network $H$; and a control coupling network $g$; propagate the state vector forward in time using the computed time derivative.

\newpage

\subsection{Innovation 2: Residual-Targeting SINDy for Satellite Manoeuvre Detection}

\textbf{Status: \textcolor{attgreen}{CLEAR — No competing patent found}}

\subsubsection*{Search Scope}
\begin{itemize}
\item Keywords: "SINDy", "sparse identification nonlinear dynamics", "satellite maneuver", "maneuver detection", "equation discovery"
\item Databases: WIPO, Google Patents, USPTO, EPO
\end{itemize}

\subsubsection*{Findings}
\begin{itemize}
\item No patents found for SINDy applied to satellite manoeuvre detection/discovery
\item SINDy base algorithm is published academic work (Brunton et al., 2016 PNAS) — \textbf{prior art in literature, not patents}
\item This means the base algorithm is freely implementable; the specific application is patentable
\end{itemize}

\subsubsection*{Novel Contribution}
\begin{itemize}
\item \textbf{Residual targeting:} $r(t) = \ddot{x}_\text{obs} - \ddot{x}_\text{Keplerian}$ rather than full dynamics $\dot{x} = f(x)$
\item \textbf{Bootstrap confidence intervals} with 1000-resample validation
\item \textbf{Hamiltonian invariant hash} for certificate issuance
\end{itemize}

\subsubsection*{Patentability Assessment}
\begin{itemize}
\item \textbf{Novelty:} High — residual-targeting approach for SSA is unpatented
\item \textbf{Non-obviousness:} Moderate — SINDy is known but residual application to Keplerian baseline is non-obvious
\item \textbf{Industrial applicability:} Clear — satellite anomaly detection for defence
\end{itemize}

\subsubsection*{Recommended IPC Classification}
\begin{itemize}
\item Primary: G06N5/00 (Computing arrangements using knowledge-based models)
\item Secondary: G05B17/02 (Automatic control systems)
\end{itemize}

\subsubsection*{Suggested Independent Claims}

\textbf{Claim 1 (Method):} A computer-implemented method for detecting satellite manoeuvres, the method comprising: obtaining observed trajectory data of a satellite; computing Keplerian-predicted trajectory data for the satellite; computing residual dynamics as a difference between observed and Keplerian-predicted accelerations; applying sparse identification of nonlinear dynamics (SINDy) to the residual dynamics to discover a sparse set of non-Keplerian terms; wherein the discovered terms represent thrust signatures indicative of a satellite manoeuvre.

\textbf{Claim 2 (Certificate Generation):} A method for generating a physics-verified manoeuvre certificate, comprising: discovering a sparse dynamical equation for satellite manoeuvre using residual-targeted SINDy; computing a Hamiltonian invariant hash from the discovered equation; generating a falsifiable prediction for a future time step; wherein the certificate is self-invalidating if Hamiltonian drift exceeds a threshold.

\newpage

\subsection{Innovation 3: IMM Filter for Hypersonic + Satellite Tracking}

\textbf{Status: \textcolor{attyellow}{PARTIALLY OVERLAPPED — Prior art exists, variant is novel}}

\subsubsection*{Search Scope}
\begin{itemize}
\item Keywords: "interacting multiple model", "IMM", "hypersonic", "satellite tracking", "space domain awareness"
\item Databases: CNIPA, Google Patents, USPTO
\end{itemize}

\subsubsection*{Prior Art Found}
\begin{itemize}
\item \textbf{CN104199022B} (Chinese, National University of Defense Technology, 2014)
  \begin{itemize}
  \item Title: "Target modal estimation based near-space hypersonic velocity target tracking method"
  \item Content: Uses IMM for near-space hypersonic tracking with dynamic mode estimation
  \item Jurisdiction: China only
  \item Status: Granted
  \end{itemize}
\end{itemize}

\subsubsection*{Overlap Analysis}
\begin{itemize}
\item \textbf{Overlapped:} General IMM framework for hypersonic tracking
\item \textbf{NOT covered by CN104199022B:}
  \begin{itemize}
  \item Bayesian online transition-matrix adaptation (Dirichlet-Markov approach)
  \item Free-energy-based adversarial intent scoring (Active Inference Expected Free Energy)
  \item \textbf{Unified multi-threat stack} — satellite manoeuvre + HGV + GNSS spoofing + debris sharing one filter architecture
  \end{itemize}
\end{itemize}

\subsubsection*{Jurisdictional Impact}
\begin{itemize}
\item Chinese patents have limited PCT/India enforceability
\item Prior art exists but is jurisdictionally and technically narrow
\end{itemize}

\subsubsection*{Patentability Assessment}
\begin{itemize}
\item \textbf{Novelty:} Moderate — base IMM is prior art, but extensions are novel
\item \textbf{Non-obviousness:} Strong — unified multi-threat architecture is non-obvious
\item \textbf{Industrial applicability:} Clear
\end{itemize}

\subsubsection*{Recommended IPC Classification}
\begin{itemize}
\item Primary: G01S13/66 (Radar tracking)
\item Secondary: G06N5/00 (Knowledge-based models)
\end{itemize}

\subsubsection*{Claim Drafting Guidance}
\begin{itemize}
\item \textbf{Avoid claiming} general IMM framework (prior art)
\item \textbf{Claim the novel extensions:}
  \begin{itemize}
  \item Unified multi-threat filtering across satellite, hypersonic, GNSS, and debris domains
  \item Bayesian online transition-matrix adaptation
  \item Free-energy-based adversarial intent scoring
  \end{itemize}
\item \textbf{Carve-around strategy:} "wherein the filter architecture is configured for simultaneous tracking across multiple threat domains"
\end{itemize}

\subsubsection*{Suggested Independent Claims}

\textbf{Claim 1 (Method):} A computer-implemented method for unified multi-threat tracking, the method comprising: maintaining a single Interacting Multiple Model (IMM) filter architecture; wherein the IMM filter is configured to simultaneously track objects across multiple threat domains including: satellite manoeuvre detection; hypersonic glide vehicle tracking; GNSS spoofing detection; and debris breakup detection; applying Bayesian online transition-matrix adaptation to estimate mode transitions; computing a free-energy-based adversarial intent score for each tracked object; wherein the unified architecture enables cross-domain threat correlation.

\newpage

\subsection{Innovation 4: Conformal Prediction for Space Domain Awareness Alerts}

\textbf{Status: \textcolor{attgreen}{CLEAR — No competing patent found}}

\subsubsection*{Search Scope}
\begin{itemize}
\item Keywords: "conformal prediction", "uncertainty quantification", "space domain awareness", "conjunction alert", "coverage certificate"
\item Databases: WIPO, Google Patents, USPTO, EPO
\end{itemize}

\subsubsection*{Findings}
\begin{itemize}
\item No patents found anywhere applying conformal prediction to SSA/conjunction alerts
\item Conformal prediction is published academic method (Vovk et al., 2005) — freely implementable
\end{itemize}

\subsubsection*{Closest Prior Art}
\begin{itemize}
\item \textbf{US6085128A} (1998, OCEAN technique) — orbital covariance estimation
  \begin{itemize}
  \item Overlap: None — classical Kalman-based, no ML, no conformal coverage guarantees
  \end{itemize}
\end{itemize}

\subsubsection*{Novel Contribution}
\begin{itemize}
\item \textbf{Distribution-free coverage certificates} on every gRPC alert
\item \textbf{Mathematically proven marginal coverage:} $P(\text{truth} \in C_\alpha) \geq 1 - \alpha$
\item Integration with sovereign air-gapped deployment
\end{itemize}

\subsubsection*{Patentability Assessment}
\begin{itemize}
\item \textbf{Novelty:} High — conformal prediction for SSA is unpatented
\item \textbf{Non-obviousness:} Strong — combining conformal ML with sovereign deployment is non-obvious
\item \textbf{Industrial applicability:} Clear
\end{itemize}

\subsubsection*{Recommended IPC Classification}
\begin{itemize}
\item Primary: G06N20/00 (Machine learning)
\item Secondary: G06Q10/04 (Forecasting or optimisation)
\end{itemize}

\subsubsection*{Suggested Independent Claims}

\textbf{Claim 1 (Method):} A computer-implemented method for generating uncertainty-quantified alerts in a space domain awareness system, the method comprising: obtaining a prediction for a space object from a machine learning model; computing a conformal prediction interval for the prediction using a calibration set; wherein the conformal prediction interval satisfies a mathematically proven marginal coverage property: $P(\text{truth} \in C_\alpha) \geq 1 - \alpha$; generating an alert comprising the prediction and the conformal prediction interval; transmitting the alert via a gRPC interface with a coverage certificate.

\newpage

\subsection{Innovation 5: NavIC Dual-Frequency ML Spoofing Detection for SSA}

\textbf{Status: \textcolor{attgreen}{LARGELY CLEAR — Generic GNSS patents exist, NavIC+SSA is novel}}

\subsubsection*{Search Scope}
\begin{itemize}
\item Keywords: "GNSS spoofing detection", "NavIC", "space situational awareness", "machine learning", "neural network"
\item Databases: EPO, USPTO, Google Patents
\end{itemize}

\subsubsection*{Prior Art Found}
\begin{itemize}
\item \textbf{EP4024087A2} — GNSS spoofing detection via bearing/range measurements
  \begin{itemize}
  \item Overlap: None — geometric/analytical approach, not ML
  \end{itemize}
\item \textbf{US20200049837A1} — ML-assisted GNSS location estimation
  \begin{itemize}
  \item Overlap: None — end-user device positioning, not SSA context
  \end{itemize}
\end{itemize}

\subsubsection*{Novel Contribution}
\begin{itemize}
\item \textbf{BiLSTM-Attention} on dual-frequency L5/S pseudorange residuals
\item \textbf{SSA operational context} — protecting a navigation constellation, not end-user positioning
\item \textbf{Dual-constellation residual architecture} specific to NavIC
\end{itemize}

\subsubsection*{Patentability Assessment}
\begin{itemize}
\item \textbf{Novelty:} High — NavIC+SSA context with ML is unpatented
\item \textbf{Non-obviousness:} Moderate — GNSS spoofing is known, but NavIC constellation protection is non-obvious
\item \textbf{Industrial applicability:} Clear
\end{itemize}

\subsubsection*{Recommended IPC Classification}
\begin{itemize}
\item Primary: G01S19/21 (GNSS signal integrity)
\item Secondary: G06N3/00 (Neural networks)
\end{itemize}

\subsubsection*{Claim Drafting Guidance}
\begin{itemize}
\item Carve around generic GNSS spoofing patents
\item Specifically claim SSA operational context
\item Claim constellation-level monitoring, not device-level
\end{itemize}

\subsubsection*{Suggested Independent Claims}

\textbf{Claim 1 (Method):} A computer-implemented method for detecting GNSS spoofing in a space domain awareness system, the method comprising: obtaining dual-frequency pseudorange measurements from a GNSS constellation; computing residuals between the dual-frequency measurements; processing the residuals using a BiLSTM-Attention neural network to classify a signal condition as normal, jamming, or spoofing; wherein the system is configured for monitoring a navigation constellation at the constellation level, not at the end-user device level.

\newpage

% =============================================================================
\section{Competitor IP Landscape}

\subsection{Commercial SSA Competitors}

\renewcommand{\arraystretch}{1.3}
\begin{tabularx}{\textwidth}{l Y c}
\toprule
\textbf{Competitor} & \textbf{Patents Found} & \textbf{Overlap Risk} \\
\midrule
LeoLabs & No ML/physics patents found & None \\
Slingshot Aerospace & No SSA ML patents; merged with Numerica SDA division (2022) & None \\
ExoAnalytic (now Anduril) & No specific ML patents; acquired by Anduril (2026) & None \\
Numerica Corp & US20140074766A1 — numerical state propagation (classical) & Low \\
Kayhan Space & No patents found in public search & None \\
\bottomrule
\end{tabularx}

\subsection{Academic/State Research}

\renewcommand{\arraystretch}{1.3}
\begin{tabularx}{\textwidth}{l Y c}
\toprule
\textbf{Entity} & \textbf{Patents Found} & \textbf{Overlap Risk} \\
\midrule
Chinese NUDT & CN104199022B — IMM hypersonic tracking & Partial on IMM \\
US DoD/AFRL & General SSA patents, no ML-physics combination & Low \\
\bottomrule
\end{tabularx}

\subsection{Geographic Jurisdiction Analysis}

\begin{itemize}
\item \textbf{United States:} Strong prior art in classical SSA methods, weak in ML-physics
\item \textbf{Europe:} Limited SSA-specific patents, strong in GNSS
\item \textbf{China:} Active in hypersonic tracking (NUDT), limited in SSA ML
\item \textbf{India:} Minimal prior art in SSA — \textbf{opportunity for first-to-file}
\end{itemize}

\newpage

% =============================================================================
\section{Filing Recommendations}

\subsection{Priority Tiering}

\textbf{🟢 GREEN — File Immediately (Highest Priority)}

\begin{enumerate}
\item \textbf{DPHNO for orbital propagation}
  \begin{itemize}
  \item No prior art
  \item Core innovation
  \item PCT application recommended
  \item IPC: G06N3/00, B64G1/10
  \end{itemize}

\item \textbf{Residual-targeting SINDy + Hamiltonian certificate}
  \begin{itemize}
  \item Unique combination
  \item No prior art
  \item PCT application recommended
  \item IPC: G06N5/00, G05B17/02
  \end{itemize}

\item \textbf{Conformal coverage certificates for conjunction alerts}
  \begin{itemize}
  \item Completely clear
  \item Novel application
  \item PCT application recommended
  \item IPC: G06N20/00, G06Q10/04
  \end{itemize}
\end{enumerate}

\textbf{🟡 YELLOW — File with Careful Claim Scoping}

\begin{enumerate}
\setcounter{enumi}{3}
\item \textbf{NavIC ML spoofing detection in SSA context}
  \begin{itemize}
  \item Carve around generic GNSS spoofing patents
  \item Specifically claim SSA operational context
  \item India + PCT filing
  \item IPC: G01S19/21, G06N3/00
  \end{itemize}

\item \textbf{Unified multi-threat IMM stack}
  \begin{itemize}
  \item Claim unified architecture, NOT base IMM
  \item India + PCT filing
  \item IPC: G01S13/66, G06N5/00
  \end{itemize}
\end{enumerate}

\subsection{Filing Strategy}

\textbf{PCT vs India:}
\begin{itemize}
\item \textbf{PCT Application:} File within 12 months of provisional, gives 30-month window for national phase entry
\item \textbf{India First:} File provisional in India first (cheaper, establishes priority date)
\item \textbf{Recommended:} File provisional in India → PCT within 12 months → national phase in US/EU/China
\end{itemize}

\textbf{Timeline (Critical):}
\begin{verbatim}
T+0:   File provisional applications in India (Section 39 secrecy direction)
T+30:  File PCT applications
T+365: Complete national phase filings (US, EP, CN)
T-7:   BEFORE IAF submission becomes public
\end{verbatim}

\textbf{Urgency:}
\begin{itemize}
\item India has \textbf{NO grace period} — public disclosure immediately bars filing
\item IAF ADITI submission will become public → file \textbf{before} submission
\item Target: File by \textbf{May 15, 2026} (10 days from now)
\end{itemize}

\newpage

\subsection{India-Specific Filing Guidance: Section 39 Secrecy Direction}

\textbf{Applicability:}
\begin{itemize}
\item Applicable to inventions "relating to atomic energy or defence"
\item DRQ/SSIS qualifies as defence-related (space domain awareness for IAF)
\end{itemize}

\textbf{Benefits:}
\begin{itemize}
\item Patent application kept secret
\item No publication until secrecy direction lifted
\item Allows disclosure to IAF without losing patent rights
\end{itemize}

\textbf{Process:}
\begin{enumerate}
\item File provisional application with IP India
\item Request Section 39 secrecy direction
\item Patent Office refers to Ministry of Defence
\item MoD grants secrecy direction (typically 6--12 months)
\item Application remains secret until direction lifted
\end{enumerate}

\textbf{Recommended Action:}
\begin{itemize}
\item File all 5 provisional applications under Section 39
\item This allows IAF submission without losing patent rights
\item Coordinate with patent attorney experienced in defence filings
\end{itemize}

\newpage

% =============================================================================
\section{Risk Assessment}

\subsection{Patent Infringement Risk}

\renewcommand{\arraystretch}{1.3}
\begin{tabularx}{\textwidth}{l Y Y}
\toprule
\textbf{Innovation} & \textbf{Infringement Risk} & \textbf{Mitigation} \\
\midrule
DPHNO & None & No prior art \\
Residual-Targeting SINDy & None & No prior art \\
IMM Filter & Low (Chinese patent) & Carve-around strategy, claim novel extensions \\
Conformal Prediction & None & No prior art \\
NavIC Spoofing & Low & Carve around generic GNSS patents, claim SSA context \\
\bottomrule
\end{tabularx}

\subsection{Freedom to Operate}

\textbf{Overall Assessment:} High — DRQ/SSIS can be deployed without significant infringement risk.

\textbf{Caveats:}
\begin{itemize}
\item Chinese IMM patent has limited enforceability outside China
\item Generic GNSS spoofing patents exist but are jurisdictionally narrow
\item No competitor holds blocking patents on ML-physics SSA approaches
\end{itemize}

\newpage

% =============================================================================
\section{Prior Art Patent Citations}

\renewcommand{\arraystretch}{1.3}
\begin{tabularx}{\textwidth}{l Y l c l}
\toprule
\textbf{Patent No.} & \textbf{Title} & \textbf{Jurisdiction} & \textbf{Year} & \textbf{Relevance} \\
\midrule
CN104199022B & Target modal estimation based near-space hypersonic velocity target tracking method & China & 2014 & Partial overlap on IMM \\
US20140074766A1 & Method and system for propagating the state of an object and its uncertainty & US & 2013 & Classical methods, no ML \\
US20230143771A1 & Space traffic management system & US (Japanese assignee) & 2023 & Data sharing architecture \\
EP4024087A2 & GNSS signal spoofing detection via bearing and/or range measurements & Europe & 2022 & Geometric approach, not ML \\
US20200049837A1 & Machine learning assisted satellite based geolocation & US & 2020 & End-user positioning, not SSA \\
US6085128A & Orbit/covariance estimation and analysis (OCEAN) & US & 1998 & Classical Kalman methods \\
\bottomrule
\end{tabularx}

\newpage

\section{Recommended Next Steps}

\subsection{Immediate (This Week)}
\begin{enumerate}
\item Engage patent attorney with experience in defence/Section 39 filings
\item Prepare provisional applications for all 5 innovations
\item File provisional applications with IP India under Section 39 secrecy direction
\end{enumerate}

\subsection{Short-Term (Next 30 Days)}
\begin{enumerate}
\item File PCT applications within 12 months of provisional
\item Coordinate with IAF on submission timeline — ensure filing before public disclosure
\item Prepare technical drawings and claim sets for national phase filings
\end{enumerate}

\subsection{Medium-Term (Next 12 Months)}
\begin{enumerate}
\item Complete national phase filings in US, EP, China
\item Monitor competitor filings — watch for blocking patents
\item File continuation applications for additional claim sets
\end{enumerate}

\vspace{2cm}

\begin{center}
\rule{0.5\textwidth}{0.4pt}\\[0.5em]
{\small\textit{Document prepared for UBERMASCHINE Pvt Ltd legal review}}\\[0.3em]
{\small\textbf{Classification: ATTORNEY-CONFIDENTIAL}}
\end{center}

\end{document}
